\chapter{Sources and Further Reading}
\label{app:sources}

Everything this book borrowed, and where to go next.

\section{The specification itself}

The normative source for everything in Part~I, and the thing to check when this
book and your memory disagree.

\begin{description}[leftmargin=1.4em, style=nextline, itemsep=4pt]
\item[The MCP specification] \texttt{modelcontextprotocol.io/specification}.
This book targets revision \texttt{2026-07-28} throughout. Older revisions stay
published, which is useful when you meet a server that speaks one of them.
\item[The schema] \texttt{schema/draft/schema.ts} in the specification
repository. TypeScript declarations with the field-level comments that explain
the defaults, and the fastest way to answer ``what exactly does this field
mean''. Several tables in this book were built by reading it directly.
\item[Extensions] \texttt{github.com/modelcontextprotocol/ext-apps} and
\texttt{ext-auth}. Extensions version independently of core, so their
repositories are the current source rather than the core specification.
\item[The SEP archive] Specification Enhancement Proposals, in
\texttt{docs/seps/}. These carry the \emph{reasoning}: SEP 2663 for the Tasks
extension and SEP 1865 for MCP Apps both explain rejected alternatives, which is
usually the part you want when a design looks arbitrary.
\end{description}

\section{The books this one is shaped by}

\begin{description}[leftmargin=1.4em, style=nextline, itemsep=4pt]
\item[\emph{High Performance Browser Networking}, Ilya Grigorik] O'Reilly, 2013,
and free at \texttt{hpbn.co}. The structural model for this book, and the source
of its title: motivate with a measurement, define the terms properly, build the
mechanism, then say what it costs. Chapter~\ref{ch:primer}'s ``speed is a
feature'' framing is lifted directly, and the latency-and-bandwidth primer is
still the clearest explanation of why round trips dominate. If you read one
thing on this list, read that one.
\item[\emph{Designing Data-Intensive Applications}, Martin Kleppmann] O'Reilly,
2017. The reference for the distributed-systems half of Part~V: idempotency,
retries, and why exactly-once delivery is a phrase to be suspicious of.
\item[\emph{Release It!}, Michael Nygard] Pragmatic Bookshelf, 2nd edition 2018.
Circuit breakers, bulkheads, and the failure patterns Chapter~\ref{ch:multi-agent}
applies to agent delegation.
\end{description}

\section{Specific ideas this book borrows}

Credit where it is owed, and where to read the original.

\begin{description}[leftmargin=1.4em, style=nextline, itemsep=4pt]
\item[The lethal trifecta] Simon Willison. Private data, untrusted content, and
an exfiltration channel; any two are survivable and all three is a breach.
Chapter~\ref{ch:threats} is organised around it because nothing else turns
prompt-injection risk into questions with yes-or-no answers.
\item[The confused deputy] Norm Hardy, ``The Confused Deputy'', ACM SIGOPS
Operating Systems Review, 1988. The original is four pages about a compiler and
a billing file, and it is worth reading precisely because the setting is so
different from ours while the problem is identical.
\item[Little's Law] John Little, 1961. $L = \lambda W$, used in
Chapter~\ref{ch:operations} to size pools. It holds for any stable system
regardless of distribution, which is why it survives contact with real traffic.
\item[Full jitter] The AWS Architecture Blog's ``Exponential Backoff and
Jitter'', which is where the specific variant in Chapter~\ref{ch:building-hosts}
comes from, along with the measurements showing why the obvious alternatives are
worse.
\end{description}

\section{The RFCs, and which parts you actually need}

Chapter~\ref{ch:authorization} profiles a stack of these. In rough order of how
often you will need to look one up:

\begin{table}[htbp]
\centering
\small
\begin{tabularx}{\textwidth}{@{}llX@{}}
\toprule
\thead{RFC} & \thead{Subject} & \thead{Read it for} \\
\midrule
7636 & PKCE                        & Why an authorization code alone is not enough \\
8707 & Resource indicators         & Audience binding, the confused-deputy fix \\
9207 & Issuer identification       & The \texttt{iss} check that stops mix-up attacks \\
9728 & Protected resource metadata & The one document your server must publish \\
8693 & Token exchange              & Delegated authorization across agents \\
6570 & URI templates               & Level 1 is all you need \\
7591 & Dynamic client registration & Deprecated; read it to migrate off it \\
\bottomrule
\end{tabularx}
\caption{The OAuth stack MCP profiles. None needs reading end to end.}
\label{tab:rfcs}
\end{table}

Also load-bearing and not an RFC: \emph{W3C Trace Context}, which defines the
\texttt{traceparent} format in Chapter~\ref{ch:building-hosts}, and \emph{JSON
Schema 2020-12}, whose composition keywords are the reason
Chapter~\ref{ch:tools} tells you to bound validation depth.

\section{Tools worth having}

\begin{description}[leftmargin=1.4em, style=nextline, itemsep=4pt]
\item[The MCP Inspector] The official interactive client. First move when a
server misbehaves and you cannot say why. Chapter~\ref{ch:testing} makes the
case for turning everything you learn in it into a contract test.
\item[The MCP Registry] \texttt{registry.modelcontextprotocol.io}. Public
discovery, and a reminder that public means untrusted.
\item[The official SDKs] TypeScript, Python, Go, Rust, Kotlin, and more. Meridian
deliberately uses none of them, because the point was to show the wire. For
anything real, use an SDK.
\end{description}

\section{On the prose}

The writing owes a debt to Steve Yegge's blog, which established that technical
writing can be long, opinionated, funny, and rigorous at the same time, and that
the reader is a colleague rather than a student.

The quality gates in \texttt{tools/} came out of trying to hold that line
mechanically: \texttt{check\_contrastive.py} exists because a first draft leaned
on ``X is not Y, it is Z'' 244 times, and \texttt{check\_listings.py} exists
because the claim that every listing is real code deserved enforcement rather
than good intentions. Appendix~\ref{ch:companion} covers running them.
